\documentclass[12pt, letterpaper]{article}

% --- Packages ---
\usepackage[margin=1in]{geometry}
\usepackage{setspace}
\usepackage{titlesec}
\usepackage{hyperref}
\usepackage{booktabs}
\usepackage{enumitem}
\usepackage{parskip}
\usepackage{fancyhdr}
\usepackage{xcolor}
\usepackage{graphicx}
\usepackage{amsmath}
\usepackage{cite}

% --- Color Definitions ---
\definecolor{titleblue}{RGB}{30, 80, 160}
\definecolor{sectiongray}{RGB}{50, 50, 50}

% --- Hyperlink Setup ---
\hypersetup{
    colorlinks=true,
    linkcolor=titleblue,
    urlcolor=titleblue,
    citecolor=titleblue
}

% --- Header / Footer ---
\pagestyle{fancy}
\fancyhf{}
\rhead{\textcolor{gray}{\small Sentiment Analysis -- Project Proposal}}
\lhead{\textcolor{gray}{\small Python Programming Course}}
\rfoot{\textcolor{gray}{\small Page \thepage}}

% --- Section Formatting ---
\titleformat{\section}{\large\bfseries\color{titleblue}}{}{0em}{}[\titlerule]
\titleformat{\subsection}{\normalsize\bfseries\color{sectiongray}}{}{0em}{}

% --- Document ---
\begin{document}

% ========================
%  TITLE PAGE
% ========================
\begin{titlepage}
    \centering
    \vspace*{2cm}

    {\Huge \textbf{\textcolor{titleblue}{Project \#11 Proposal}}}\\[0.5cm]
    {\LARGE Reddit \& Twitter Sentiment Analysis}\\[0.3cm]
    {\large A Computational Social Science Study}

    \vspace{1.5cm}
    \rule{0.6\textwidth}{0.4pt}
    \vspace{1.5cm}

    {\large \textbf{Course:} Python Programming}\\[0.3cm]
    {\large \textbf{Student:} Vikram}\\[0.3cm]
    {\large \textbf{Date:} \today}

    \vfill
    \textcolor{gray}{\small Submitted as a project proposal for approval}
\end{titlepage}

% ========================
%  TABLE OF CONTENTS
% ========================
\tableofcontents
\newpage

% ========================
%  1. INTRODUCTION
% ========================
\section{Introduction}

Social media platforms like Reddit and Twitter (now X) generate millions of posts daily, making them rich, real-time repositories of public opinion. Sentiment analysis---the computational study of emotional tone in text---offers a powerful lens for understanding how communities feel about events, brands, products, or social issues over time.

This project proposes building a complete data pipeline in Python that collects social media data, performs sentiment analysis, and produces meaningful visualizations to uncover trends in public discourse.

\vspace{0.5em}
\noindent \textbf{Core Question:} \textit{How does public sentiment on a chosen topic (e.g., AI, mental health, climate) evolve over time across Reddit communities, and what linguistic patterns characterize positive vs. negative discourse?}

% ========================
%  2. DATA
% ========================
\section{Data Sources}

\subsection{Primary Dataset}

Two options will be evaluated; the final choice depends on API access at time of execution:

\begin{table}[h!]
\centering
\renewcommand{\arraystretch}{1.4}
\begin{tabular}{@{}lll@{}}
\toprule
\textbf{Source} & \textbf{Access Method} & \textbf{Notes} \\
\midrule
Reddit & PRAW (Python Reddit API Wrapper) & Free, well-documented \\
Twitter / X & Tweepy or Academic API & Rate-limited on free tier \\
Kaggle Datasets & Direct CSV download & Pre-scraped, no API needed \\
\bottomrule
\end{tabular}
\caption{Candidate data sources for the project}
\end{table}

\subsection{Preferred Approach}

The project will use the \textbf{Reddit API via PRAW} as the primary source, targeting 2--3 subreddits on a single topic (e.g., \texttt{r/technology}, \texttt{r/artificialintelligence}, \texttt{r/MachineLearning}).

A pre-scraped \textbf{Kaggle dataset} will serve as a backup if API limits are encountered.

\subsection{Data Fields Collected}

\begin{itemize}[noitemsep]
    \item Post title and body text
    \item Subreddit name
    \item Post creation timestamp
    \item Upvote score and ratio
    \item Number of comments
    \item Author name (anonymized for reporting)
\end{itemize}

\subsection{Estimated Volume}

Approximately 5,000--15,000 posts will be collected, spanning a 6--12 month window to enable meaningful time-series analysis.

% ========================
%  3. ANALYSIS PLAN
% ========================
\section{Proposed Analysis}

\subsection{Data Cleaning \& Preprocessing}

\begin{enumerate}[noitemsep]
    \item Remove deleted/removed posts and bot accounts
    \item Strip URLs, special characters, and HTML entities
    \item Tokenize and lowercase all text
    \item Remove stopwords using \texttt{NLTK}
    \item Apply lemmatization for word normalization
\end{enumerate}

\subsection{Sentiment Scoring}

Sentiment will be assigned using two complementary methods:

\begin{itemize}
    \item \textbf{VADER} (\textit{Valence Aware Dictionary and sEntiment Reasoner}): A lexicon-based model well-suited for social media text. Outputs a compound score in $[-1, +1]$.
    \item \textbf{TextBlob}: A simpler alternative providing polarity and subjectivity scores; useful for cross-validation.
\end{itemize}

Each post will be classified as \textit{Positive}, \textit{Neutral}, or \textit{Negative} based on the VADER compound score:
\[
\text{Label} =
\begin{cases}
\text{Positive} & \text{if } c \geq 0.05 \\
\text{Negative} & \text{if } c \leq -0.05 \\
\text{Neutral} & \text{otherwise}
\end{cases}
\]

\subsection{Statistical Analyses}

\begin{enumerate}
    \item \textbf{Sentiment over time:} Compute weekly rolling averages of compound sentiment scores.
    \item \textbf{Engagement vs.\ sentiment:} Pearson correlation between upvote score and sentiment polarity.
    \item \textbf{Subreddit comparison:} ANOVA or Mann-Whitney U test to compare mean sentiment across subreddits.
    \item \textbf{Top keywords by sentiment class:} TF-IDF to extract the most discriminating words per class.
\end{enumerate}

% ========================
%  4. VISUALIZATION PLAN
% ========================
\section{Proposed Visualizations}

All visualizations will be produced using \texttt{matplotlib}, \texttt{seaborn}, and \texttt{wordcloud} in Python.

\begin{table}[h!]
\centering
\renewcommand{\arraystretch}{1.4}
\begin{tabular}{@{}p{3.5cm}p{4.5cm}p{5cm}@{}}
\toprule
\textbf{Chart Type} & \textbf{Variables} & \textbf{Purpose} \\
\midrule
Line Chart & Sentiment score vs.\ time & Show mood trends over weeks/months \\
Bar Chart & Avg.\ sentiment by subreddit & Compare community tone \\
Scatter Plot & Upvote score vs.\ sentiment & Explore engagement--emotion link \\
Word Cloud & Top words per sentiment class & Visual summary of vocabulary \\
Histogram & Distribution of compound scores & Understand score spread \\
Heatmap & Sentiment by day-of-week \& hour & Temporal posting patterns \\
\bottomrule
\end{tabular}
\caption{Planned visualizations and their analytical purpose}
\end{table}

% ========================
%  5. TOOLS & LIBRARIES
% ========================
\section{Tools and Libraries}

\begin{itemize}[noitemsep]
    \item \textbf{Data Collection:} \texttt{praw}, \texttt{requests}, \texttt{pandas}
    \item \textbf{Text Processing:} \texttt{nltk}, \texttt{re}, \texttt{string}
    \item \textbf{Sentiment Analysis:} \texttt{vaderSentiment}, \texttt{textblob}
    \item \textbf{Visualization:} \texttt{matplotlib}, \texttt{seaborn}, \texttt{wordcloud}
    \item \textbf{Statistics:} \texttt{scipy.stats}, \texttt{sklearn} (TF-IDF)
    \item \textbf{Report:} \LaTeX{} (\texttt{pdflatex})
\end{itemize}

% ========================
%  6. PROJECT TIMELINE
% ========================
\section{Project Timeline}

\begin{table}[h!]
\centering
\renewcommand{\arraystretch}{1.4}
\begin{tabular}{@{}lll@{}}
\toprule
\textbf{Week} & \textbf{Task} & \textbf{Deliverable} \\
\midrule
Week 1 & Data collection \& cleaning & Clean CSV dataset \\
Week 2 & Sentiment scoring \& EDA & Annotated DataFrame \\
Week 3 & Statistical analysis & Analysis notebook \\
Week 4 & Visualizations & Publication-ready figures \\
Week 5 & LaTeX report writing & Final PDF report \\
\bottomrule
\end{tabular}
\caption{Proposed five-week project schedule}
\end{table}

% ========================
%  7. EXPECTED OUTCOMES
% ========================
\section{Expected Outcomes}

Upon completion, this project will produce:

\begin{enumerate}
    \item A clean, reusable Python data pipeline for social media NLP tasks
    \item A statistical analysis revealing sentiment trends and engagement patterns
    \item Six or more original visualizations suitable for a research report
    \item A full LaTeX report documenting methodology, findings, and conclusions
\end{enumerate}

The project will demonstrate proficiency in data collection, text preprocessing, applied NLP, statistical thinking, and scientific communication---core skills for data science and computational research roles.

% ========================
%  8. CONCLUSION
% ========================
\section{Conclusion}

Reddit and Twitter offer an unparalleled window into real-time public sentiment. By applying Python-based NLP tools to a focused corpus of posts, this project will yield actionable insights into how communities discuss a topic of interest---and how emotional tone correlates with online engagement. The methods are scalable, reproducible, and directly applicable to real-world problems in social science, public health communication, and market research.

% ========================
%  REFERENCES
% ========================
\newpage
\section*{References}

\begin{enumerate}
    \item Hutto, C.J. \& Gilbert, E. (2014). VADER: A Parsimonious Rule-based Model for Sentiment Analysis of Social Media Text. \textit{ICWSM}.
    \item Bonta, V. \& Janardhan, N. (2019). A Comprehensive Study on Lexicon-Based Approaches for Sentiment Analysis. \textit{Asian Journal of Computer Science and Technology}.
    \item Reddit API Documentation. \url{https://www.reddit.com/dev/api/}
    \item PRAW Documentation. \url{https://praw.readthedocs.io/}
    \item Kaggle Datasets Portal. \url{https://www.kaggle.com/datasets}
\end{enumerate}

\end{document}
